% =============================================================================
% CAPÍTULO 1 — INTRODUCCIÓN Y ASPECTOS GENERALES (ISO/IEC/IEEE 15289:2024)
% =============================================================================
\chapter{Introducción y Aspectos Generales}
\label{ch:introduccion}
\resumenCapitulo{Este capítulo presenta el propósito, el alcance y la audiencia del manual,
fija el glosario y las referencias normativas, y describe la política de versionado del
documento y de la plataforma. Constituye el marco de referencia que contextualiza los
capítulos técnicos siguientes.}

Este capítulo establece el marco de referencia del presente Manual Técnico. Conforme
al estándar \textbf{ISO/IEC/IEEE 15289:2024} —que define el contenido de los productos
de información del ciclo de vida del software y de los sistemas— se delimita el
propósito, la audiencia, el alcance, el vocabulario controlado y la estrategia de
trazabilidad que rigen el resto del documento.

El estándar ISO/IEC/IEEE 15289 clasifica los productos de información en cuatro tipos:
\textit{descripción} (\textit{description}), \textit{plan}, \textit{política}
(\textit{policy}), \textit{procedimiento} (\textit{procedure}), \textit{reporte},
\textit{solicitud}, \textit{especificación} (\textit{specification}) y otros. El
presente Manual Técnico constituye principalmente una \textbf{descripción de diseño
arquitectónico} (\textit{architecture description}) y una \textbf{especificación de
diseño} (\textit{design specification}), complementada con elementos de
especificación de requisitos regidos por \textbf{IEEE 830}. En consecuencia, todo su
contenido se subordina a un principio rector: \textit{ser verificable contra el código
fuente real de la plataforma}, evitando descripciones genéricas o aspiracionales.

\begin{AvisoNota}
Cada afirmación técnica de este manual es trazable a un artefacto concreto del
repositorio: un paquete Java, un componente React, una migración SQL o un archivo de
configuración. Las referencias a clases, métodos y rutas emplean tipografía
\comando{monoespaciada} precisamente para resaltar dicha trazabilidad.
\end{AvisoNota}

\paragraph{Organización del documento.} El manual se estructura en nueve capítulos que
recorren el sistema desde lo general a lo particular: la introducción y los requisitos
(capítulos~\ref{ch:introduccion} y~\ref{ch:requisitos}) fijan el \textit{qué}; la
arquitectura global C4 (capítulo~\ref{ch:arquitectura}) describe el \textit{cómo} a
alto nivel; los capítulos de diseño detallado de backend, frontend y base de datos
descienden al \textit{cómo} concreto; y los capítulos finales cubren las integraciones
externas, la verificación y validación, y los anexos documentales.

% -----------------------------------------------------------------------------
\section{Propósito y Audiencia del Documento}
\label{sec:proposito}

El propósito de este manual es servir como \textbf{fuente única y autoritativa} de la
arquitectura, el diseño detallado y las decisiones técnicas de la plataforma
\textbf{EthosHub} en su versión \texttt{1.1.0}. A diferencia de un manual de usuario,
este documento describe el \textit{cómo} y el \textit{porqué} de la construcción del
sistema, habilitando su mantenimiento, evolución y operación a lo largo del tiempo.

El documento está dirigido a tres perfiles profesionales:

\begin{itemize}[noitemsep]
    \item \textbf{Ingenieros de software (Backend y Frontend):} responsables de
    mantener y evolucionar el monolito Spring Boot y la SPA React. Encontrarán
    diagramas de clases, secuencias de operaciones clave, contratos de la API REST y
    la organización del código por \textit{feature}.
    \item \textbf{Administradores de bases de datos (DBA):} responsables del esquema
    \texttt{core} sobre PostgreSQL/Supabase. Encontrarán el diccionario de datos
    exhaustivo, las políticas de integridad referencial, la estrategia de indexación
    y las políticas de seguridad a nivel de fila (RLS).
    \item \textbf{Personal de operaciones (DevOps):} responsable del empaquetado, el
    despliegue y la observabilidad. Encontrarán el flujo de construcción monolítica
    (\comando{build.sh}), la especificación del contenedor Docker y la gestión de
    configuraciones y secretos.
\end{itemize}

Dado que cada perfil consulta el manual con objetivos distintos, la
tabla~\ref{tab:guia-lectura} ofrece una \textbf{guía de lectura orientada por rol},
indicando los capítulos de mayor relevancia para cada audiencia y los conocimientos
previos que se asumen.

\vspace{0.2cm}
\noindent
\renewcommand{\arraystretch}{1.4}
\fontsize{9}{10.8}\selectfont\sffamily
\begin{tabularx}{\textwidth}{|>{\bfseries\color{colorPrimario}}l|X|c|}
\hline
\rowcolor{headerTabla}
\textbf{Perfil} & \textbf{Conocimientos asumidos y capítulos prioritarios} & \textbf{Capítulos} \\
\hline
Ingeniero Backend & Java 17, Spring Boot, OAuth2, SQL tipado (jOOQ). Diseño de la API, secuencias y excepciones. & 3, 4, 6, 7 \\
\hline
\rowcolor{grisTecnico}
Ingeniero Frontend & React 18, TypeScript estricto, Zustand, React Router. Componentes, \textit{stores}, interceptores. & 3, 5, 7 \\
\hline
DBA & PostgreSQL, RLS, indexación, PL/pgSQL. Diccionario de datos, integridad y triggers. & 6 \\
\hline
\rowcolor{grisTecnico}
DevOps / Operaciones & Docker, variables de entorno, CI. Empaquetado, entornos y gestión de secretos. & 3, 8, 9 \\
\hline
\end{tabularx}
\label{tab:guia-lectura}

\begin{AvisoNota}
Para tareas de instalación y puesta en marcha paso a paso consulte el \textit{Manual
de Instalación}; para la operativa funcional de la plataforma, el \textit{Manual de
Usuario}. Este manual asume conocimientos sólidos de Java, Spring, React, SQL y
contenedores, por lo que no explica dichas tecnologías base, sino su \textbf{aplicación
concreta} en EthosHub.
\end{AvisoNota}

% -----------------------------------------------------------------------------
\section{Alcance del Sistema}
\label{sec:alcance}

\textbf{EthosHub} es una plataforma unificada de \textit{Talent Hub} que articula dos
experiencias complementarias sobre un mismo dominio de datos:

\begin{enumerate}[noitemsep]
    \item \textbf{Construcción de portafolios profesionales.} El profesional
    construye un portafolio público enriquecido: proyectos (con documentos PDF y
    archivos importados desde Google Drive), experiencia laboral, formación
    académica, habilidades técnicas y blandas, y un currículum interactivo generado
    en el \textit{CV Studio} con asistencia de IA.
    \item \textbf{Descubrimiento de talento.} El reclutador descubre, filtra y conecta
    con talento mediante un \textit{pipeline} tipo Kanban, métricas, un
    \textit{leaderboard} e \textit{insights} asistidos por inteligencia artificial.
\end{enumerate}

El alcance técnico documentado cubre el \textbf{monolito empaquetado} que integra el
backend (\texttt{ethoshubB}, Spring Boot) y el frontend (\texttt{ethoshubF}, React SPA)
en un único artefacto ejecutable, junto con sus dependencias externas: \textbf{Supabase}
(autenticación, base de datos PostgreSQL y comunicación en tiempo real),
\textbf{Google Gemini} (servicios de IA) y \textbf{Google Drive} (importación de
archivos).

El \textbf{stack tecnológico} que materializa este alcance se resume en la
tabla~\ref{tab:stack-global}, organizado por capa. Estas elecciones tecnológicas
condicionan transversalmente el diseño documentado en los capítulos siguientes.

\clearpage
\vspace{0.2cm}
\noindent
\renewcommand{\arraystretch}{1.35}
\fontsize{8.5}{10.2}\selectfont\sffamily
\begin{tabularx}{\textwidth}{|>{\bfseries\color{colorPrimario}}l|l|X|}
\hline
\rowcolor{headerTabla}
\textbf{Capa} & \textbf{Tecnología} & \textbf{Rol en EthosHub} \\
\hline
Lenguaje (BE) & Java 17 & Lenguaje del monolito Spring Boot. \\
\hline
\rowcolor{grisTecnico}
Framework (BE) & Spring Boot 4.0.x & Web MVC, Security (Resource Server), Validation, Actuator, Mail. \\
\hline
Persistencia & jOOQ + Spring JDBC & SQL tipado contra el esquema \texttt{core}. \\
\hline
\rowcolor{grisTecnico}
Base de datos & PostgreSQL (Supabase) & Almacenamiento relacional con RLS y RPCs \texttt{SECURITY DEFINER}. \\
\hline
Lenguaje (FE) & TypeScript 5.3 (\textit{strict}) & Tipado estático de la SPA. \\
\hline
\rowcolor{grisTecnico}
Framework (FE) & React 18.3 + Vite 5.1 & SPA y servidor de desarrollo con HMR. \\
\hline
Estado (FE) & Zustand 4.5 & \textit{Stores} por dominio con persistencia en \textit{sessionStorage}. \\
\hline
\rowcolor{grisTecnico}
IA & Google Gemini (REST) & Asistente de CV Studio e \textit{insights} de reclutamiento. \\
\hline
Documentación API & springdoc-openapi 2.7 & Generación de Swagger UI y OpenAPI. \\
\hline
\rowcolor{grisTecnico}
Empaquetado & Maven + Docker & JAR monolítico y contenedor con Tectonic. \\
\hline
\end{tabularx}
\label{tab:stack-global}

\vspace{0.4cm}
Para delimitar con precisión qué responsabilidades asume EthosHub y cuáles delega, la
figura~\ref{fig:alcance} contrasta los elementos \textbf{dentro} y \textbf{fuera} del
alcance del sistema.

\vspace{0.3cm}
\begin{figure}[h!]
\centering
\begin{tikzpicture}[
    item/.style={rectangle, rounded corners=2pt, draw, font=\sffamily\scriptsize, align=center,
        minimum width=3.4cm, minimum height=0.7cm, inner sep=3pt},
    dentro/.style={item, fill=c4Componente!30, draw=c4Sistema!70!black},
    fuera/.style={item, fill=c4Externo!20, draw=c4Externo!70!black, dashed}
]
    \node[font=\sffamily\footnotesize\bfseries\color{c4Sistema}] (td) at (-3.2,4.1) {Dentro del alcance};
    \node[font=\sffamily\footnotesize\bfseries\color{c4Externo!60!black}] (tf) at (3.2,4.1) {Fuera del alcance (delegado)};

    \node[dentro] (d1) at (-3.2,3.3) {Portafolios y contenido profesional};
    \node[dentro] (d2) at (-3.2,2.4) {Descubrimiento de talento};
    \node[dentro] (d3) at (-3.2,1.5) {Empaquetado monolítico};
    \node[dentro] (d4) at (-3.2,0.6) {Validación de JWT (Resource Server)};
    \node[dentro] (d5) at (-3.2,-0.3) {Orquestación de IA y exportación de CV};

    \node[fuera] (f1) at (3.2,3.3) {Emisión de credenciales (Supabase Auth)};
    \node[fuera] (f2) at (3.2,2.4) {Hosting de la base de datos (Supabase)};
    \node[fuera] (f3) at (3.2,1.5) {Modelos de IA (Google Gemini)};
    \node[fuera] (f4) at (3.2,0.6) {Almacenamiento de archivos (Drive/Supabase)};
    \node[fuera] (f5) at (3.2,-0.3) {Entrega de correo (SMTP Gmail)};

    \draw[grisISO, thick, dashed] (0,4.0) -- (0,-0.7);
\end{tikzpicture}
\caption{Delimitación del alcance: responsabilidades propias frente a servicios delegados.}
\label{fig:alcance}
\end{figure}

\begin{NotaTecnica}
\textbf{Fuera de alcance.} La autenticación de usuarios se delega íntegramente en
\textbf{Supabase Auth}: el backend actúa como \textit{Resource Server} OAuth2 y solo
valida los \textit{tokens} JWT emitidos por Supabase; no implementa un servidor de
autorización propio. Esta decisión reduce la superficie de ataque del backend y
externaliza la gestión de contraseñas, MFA y proveedores OAuth a un servicio
especializado.
\end{NotaTecnica}

% -----------------------------------------------------------------------------
\section{Definiciones, Acrónimos y Abreviaturas}
\label{sec:definiciones}

\subsection{Glosario de términos arquitectónicos}
\label{subsec:glosario_arq}

\vspace{0.2cm}
\noindent
\renewcommand{\arraystretch}{1.4}
\fontsize{9}{10.8}\selectfont\sffamily
\begin{tabularx}{\textwidth}{|>{\bfseries\color{colorPrimario}}l|X|}
\hline
\rowcolor{headerTabla}
\textbf{Término} & \textbf{Definición} \\
\hline
BaaS & \textit{Backend as a Service}. Modelo en el que servicios de infraestructura (autenticación, base de datos, tiempo real, almacenamiento) se consumen como servicio gestionado. En EthosHub, este rol lo cumple \textbf{Supabase}. \\
\hline
\rowcolor{grisTecnico}
SPA & \textit{Single Page Application}. Aplicación web que carga una sola página HTML y reescribe dinámicamente su contenido vía JavaScript. El frontend de EthosHub es una SPA React. \\
\hline
Monolito Empaquetado & Patrón de despliegue en el que la SPA compilada (\comando{dist/}) se copia al directorio estático de Spring Boot y se sirve desde el mismo JAR ejecutable que la API REST. Un único artefacto despliega API y cliente. \\
\hline
\rowcolor{grisTecnico}
Modelo C4 & Notación de \textit{Context, Containers, Components, Code} para describir arquitecturas de software por niveles de abstracción. Se emplea en el capítulo~\ref{ch:arquitectura}. \\
\hline
\textit{Vertical Slice} & Organización del código por \textit{feature} (dominio funcional) en lugar de por capa técnica: cada paquete agrupa su \textit{controller}, \textit{service}, \textit{repository} y DTOs. \\
\hline
\rowcolor{grisTecnico}
Resource Server & Componente OAuth2 que protege recursos validando \textit{tokens} de acceso (JWT) sin emitirlos. \\
\hline
RLS & \textit{Row Level Security}. Mecanismo de PostgreSQL que restringe el acceso a filas individuales según políticas asociadas al rol del usuario. \\
\hline
\rowcolor{grisTecnico}
RPC \texttt{SECURITY DEFINER} & Función de PostgreSQL que se ejecuta con los privilegios de su propietario, usada para operar de forma controlada por encima de las políticas RLS. \\
\hline
JWT / JWKS & \textit{JSON Web Token} y \textit{JSON Web Key Set}. El JWT transporta la identidad firmada; el JWKS publica las claves públicas para validar dicha firma. \\
\hline
\end{tabularx}

\subsection{Definiciones de dominio de negocio}
\label{subsec:glosario_negocio}

\vspace{0.2cm}
\noindent
\renewcommand{\arraystretch}{1.4}
\fontsize{9}{10.8}\selectfont\sffamily
\begin{tabularx}{\textwidth}{|>{\bfseries\color{colorPrimario}}l|X|}
\hline
\rowcolor{headerTabla}
\textbf{Término} & \textbf{Definición} \\
\hline
Perfil & Identidad profesional de un usuario en la plataforma. Es la entidad raíz del dominio (tabla \comando{core.profiles}); de ella cuelgan proyectos, experiencia, formación y habilidades. \\
\hline
\rowcolor{grisTecnico}
Portafolio & Vista pública configurable del Perfil: agrega contenido profesional con ajustes de visibilidad, SEO y, opcionalmente, protección por contraseña. \\
\hline
CV Studio & Editor avanzado de currículums basado en Monaco Editor, con asistente de IA Gemini multi-turno y compilación a PDF mediante LaTeX (Tectonic). \\
\hline
\rowcolor{grisTecnico}
Espacio de Reclutamiento & Conjunto de funcionalidades del Perfil Reclutador: descubrimiento de talento, \textit{pipeline} Kanban de candidatos, notas privadas, \textit{leaderboard} e \textit{insights} de IA. \\
\hline
\end{tabularx}

% -----------------------------------------------------------------------------
\section{Referencias y Estándares Aplicados}
\label{sec:referencias}

La elaboración de este manual se rige por el cumplimiento normativo de los siguientes
estándares de ingeniería:

\begin{itemize}[noitemsep]
    \item \textbf{ISO/IEC/IEEE 15289:2024} — \textit{Systems and software engineering
    — Content of life-cycle information items (documentation)}. Define la estructura,
    el contenido mínimo y la trazabilidad de los productos de información, incluyendo
    el historial de revisiones, los roles y la identificación documental.
    \item \textbf{IEEE 830} — \textit{Recommended Practice for Software Requirements
    Specifications}. Rige la especificación de requisitos del capítulo~\ref{ch:requisitos}:
    actores, casos de uso, requisitos funcionales por dominio y requisitos no
    funcionales (rendimiento, concurrencia, disponibilidad).
    \item \textbf{Modelo C4 (Simon Brown)} — Notación complementaria empleada para la
    descripción de la arquitectura por niveles (contexto, contenedores y componentes).
\end{itemize}

Los documentos contractuales y de origen de la plataforma se listan en la bibliografía:
la convocatoria pública CPTIS-2302-2026 y su pliego de especificaciones constituyen el
marco de requisitos originales del sistema.

Para evidenciar el \textbf{cumplimiento normativo}, la tabla~\ref{tab:cumplimiento}
mapea las cláusulas relevantes de cada estándar con el capítulo del manual que las
satisface. Este mapeo permite a un evaluador verificar, de forma directa, que cada
exigencia del estándar tiene un correlato concreto en el documento.

\vspace{0.2cm}
\noindent
\renewcommand{\arraystretch}{1.35}
\fontsize{9}{10.8}\selectfont\sffamily
\begin{tabularx}{\textwidth}{|>{\bfseries\color{colorPrimario}}l|X|c|}
\hline
\rowcolor{headerTabla}
\textbf{Estándar} & \textbf{Exigencia / producto de información} & \textbf{Capítulo} \\
\hline
ISO 15289 & Identificación y control de versiones del documento & Historial \\
\hline
\rowcolor{grisTecnico}
ISO 15289 & Roles y responsabilidades sobre el artefacto & Historial \\
\hline
ISO 15289 & Descripción de la arquitectura del sistema & 3 \\
\hline
\rowcolor{grisTecnico}
ISO 15289 & Especificación de diseño detallado & 4, 5, 6 \\
\hline
ISO 15289 & Información de verificación y validación & 8 \\
\hline
\rowcolor{grisTecnico}
IEEE 830 & Actores y casos de uso del sistema & 2 \\
\hline
IEEE 830 & Requisitos funcionales específicos & 2 \\
\hline
\rowcolor{grisTecnico}
IEEE 830 & Requisitos no funcionales (rendimiento, restricciones) & 2 \\
\hline
Modelo C4 & Vistas de contexto, contenedores y componentes & 3 \\
\hline
\end{tabularx}
\label{tab:cumplimiento}

% -----------------------------------------------------------------------------
\section{Estrategia de Versionado y Trazabilidad}
\label{sec:versionado}

EthosHub adopta \textbf{Versionado Semántico} (\textit{SemVer},
\texttt{MAYOR.MENOR.PARCHE}) tanto para el producto como para sus artefactos
documentales. La semántica de cada segmento se define en la
tabla~\ref{tab:semver}, y rige tanto la numeración del JAR
(\texttt{ethoshub-0.0.1-SNAPSHOT}) como la de este manual (\texttt{1.1.0}).

\vspace{0.2cm}
\noindent
\renewcommand{\arraystretch}{1.35}
\fontsize{9}{10.8}\selectfont\sffamily
\begin{tabularx}{\textwidth}{|>{\bfseries\color{colorPrimario}}l|X|l|}
\hline
\rowcolor{headerTabla}
\textbf{Segmento} & \textbf{Criterio de incremento} & \textbf{Ejemplo} \\
\hline
MAYOR & Cambios incompatibles en la API o el modelo de datos (rupturas de contrato). & 1.x $\rightarrow$ 2.0.0 \\
\hline
\rowcolor{grisTecnico}
MENOR & Nueva funcionalidad retrocompatible (nuevo dominio o \textit{endpoint}). & 2.0.0 $\rightarrow$ 1.1.0 \\
\hline
PARCHE & Correcciones retrocompatibles (\textit{bug fixes}, ajustes de RLS). & 1.1.0 $\rightarrow$ 1.1.1 \\
\hline
\end{tabularx}
\label{tab:semver}

\vspace{0.4cm}
La trazabilidad se garantiza mapeando cada incremento entre tres planos: los
\textbf{Sprints} de desarrollo, las \textbf{Historias de Perfil} (\textit{User
Stories}) y los \textbf{componentes de software} afectados. El siguiente diagrama
resume dicho mapeo.

\vspace{0.3cm}
\begin{figure}[h!]
\centering
\begin{tikzpicture}[node distance=0.9cm and 1.6cm]
    % Capas
    \node[c4persona, fill=c4Persona, minimum width=3.2cm] (sprint) {Sprint\\\scriptsize(iteración de desarrollo)};
    \node[c4contenedor, right=of sprint, minimum width=3.2cm] (us) {Historia de Perfil\\\scriptsize(User Story)};
    \node[c4componente, right=of us, minimum width=3.2cm] (comp) {Componente de software\\\scriptsize(paquete / módulo)};

    \draw[c4flecha] (sprint) -- node[c4etiqueta, above]{planifica} (us);
    \draw[c4flecha] (us) -- node[c4etiqueta, above]{implementa} (comp);

    % Ejemplos concretos
    \node[draw=grisISO!50, rounded corners, fill=grisTecnico, font=\sffamily\scriptsize, align=center, below=1.1cm of sprint, minimum width=3.2cm] (e1) {S1 — Identidad y Auth};
    \node[draw=grisISO!50, rounded corners, fill=grisTecnico, font=\sffamily\scriptsize, align=center, below=1.1cm of us, minimum width=3.2cm] (e2) {LOG-001 Registro\\IDEN-003 Biografía};
    \node[draw=grisISO!50, rounded corners, fill=grisTecnico, font=\sffamily\scriptsize, align=center, below=1.1cm of comp, minimum width=3.2cm] (e3) {\texttt{auth/}, \texttt{profile/}\\\texttt{authStore.ts}};

    \draw[c4flecha, dashed] (sprint) -- (e1);
    \draw[c4flecha, dashed] (us) -- (e2);
    \draw[c4flecha, dashed] (comp) -- (e3);
\end{tikzpicture}
\caption{Cadena de trazabilidad entre Sprints, Historias de Perfil y componentes de software.}
\label{fig:trazabilidad}
\end{figure}

Esta cadena permite responder, ante cualquier cambio, a las preguntas de auditoría
fundamentales: \textit{¿en qué iteración se introdujo?}, \textit{¿qué necesidad de
negocio lo motivó?} y \textit{¿qué artefactos de código lo materializan?}. La
tabla~\ref{tab:migraciones-trazabilidad} ilustra este principio sobre las migraciones
de base de datos, donde cada script \comando{V$n$} es trazable a un incremento
funcional.

\vspace{0.2cm}
\noindent
\renewcommand{\arraystretch}{1.35}
\fontsize{9}{10.8}\selectfont\sffamily
\begin{tabularx}{\textwidth}{|c|X|l|}
\hline
\rowcolor{headerTabla}
\textbf{Migración} & \textbf{Incremento funcional} & \textbf{Dominio} \\
\hline
V1 & Campos geográficos en experiencias laborales & \texttt{experience} \\
\hline
\rowcolor{grisTecnico}
V6 & Fases de \textit{Talent Discovery} (Kanban, notas, dashboard) & \texttt{recruiter} \\
\hline
V7 & Idempotencia de mensajes y RPCs de notas seguras & \texttt{chat} \\
\hline
\rowcolor{grisTecnico}
V14 & Conexiones de perfil con enfoque \textit{URL-first} (+15 \textit{providers}) & \texttt{connections} \\
\hline
V15 & Política de SELECT para \textit{likes} de perfil & \texttt{recruiter} \\
\hline
\end{tabularx}
\label{tab:migraciones-trazabilidad}
